\section{Project Plan}
%=====================================================================

\subsection{Activities and Process}

\revised{The team follows an iterative process in six phases of one to two weeks each (see the timetable). Each phase ends with a deliverable that someone can check, and the plan lives in a GitHub issue that every member can read. Two milestones, Midterm report (Nov 9) and Final report (Dec 11), group the issues of the phases that end before each date.}

Every piece of work is one GitHub issue. The issue belongs to a milestone and to the Capstone project board, which holds its start and target dates. An issue is sized for one person and about one week: larger work is split, and work under a day joins a related issue. Each issue ends with a ``Done when'' list that a second person can check, and one line of requirement IDs so the work traces back to the SRS.

Each issue gets its own branch, and nothing is committed to \texttt{main} directly. Changes enter through one pull request that links its issue. At least one member besides the author reviews the pull request before it merges. Documents follow the same path as code.

Members pick the issues they want, except setup issues that other work depends on, which go first. After the meetings with CTI and with the Athletic Activities department, the team held standups where decisions were made. Meeting summaries and decisions are posted as announcements in the repository's Discussions with their dates. \revised{When a decision changes the plan, the team updates the plan issue, the affected issues, and the Testing Plan.}

\subsection{Timetable, Milestones, and Deliverables}

\begin{center}
\small
\begin{longtable}{L{1.3cm} L{2.4cm} L{11.3cm}}
\toprule
\textbf{Phase} & \textbf{Dates} & \textbf{Deliverable at the end of the phase} \\
\midrule
\endfirsthead
\toprule
\textbf{Phase} & \textbf{Dates} & \textbf{Deliverable at the end of the phase} \\
\midrule
\endhead
\bottomrule
\endlastfoot
1 & Oct 14 to Oct 25 & Decisions and screens: the team agreed on every open design decision and every screen has an approved wireframe. \\
\addlinespace
2 & Oct 19 to Nov 1 & Foundation: one command starts everything, the PWA opens on a phone, the database exists, and a person can sign in in Spanish or English. \\
\addlinespace
3 & Nov 2 to Nov 8 & Check In: a person with an authorization does Check In inside the gym, sees the Gym ID, and sees the estimated occupancy. An Administrator creates terms and the athlete roster. \textbf{Midterm report, Nov 9.} \\
\addlinespace
4 & Nov 9 to Nov 22 & Registration and sync: people register as paid members or athletes, renew, and become athletes inside their account, and the app reads paid orders from the store. \\
\addlinespace
5 & Nov 16 to Nov 29 & Staff and administration: Staff see active sessions, verify Gym IDs, and do assisted Check In; Administrators manage access, schedules, and notices; everyone sees analytics. \\
\addlinespace
6 & Nov 30 to Dec 11 & Release: the app runs on the CTI server and the gym accepted it. \textbf{Final report, Dec 11.} \\
\end{longtable}
\end{center}

\subsection{Team Responsibilities}

\revised{Each member owns work areas and reviews the others' pull requests. All three met with CTI and the gym and wrote this report. Table~\ref{tab:team} lists what each member did and owns.}

\begin{table}[h]
\centering
\caption{\revised{Work done and owned by each member}}
\label{tab:team}
\small
\begin{tabularx}{\textwidth}{L{3.2cm} X}
\toprule
\textbf{Member} & \textbf{Work done and owned} \\
\midrule
\revised{Carlos J. Pepin Delgado} & \revised{Software Requirements Specification and design documents: architecture, data model, registration, Check In and administration flows, and the screen map by role. Team rules, issue templates, and the wireframe guide. Frontend base and its CI. User app base: phone layout, navigation, install help, and offline support. User sign in on one device, password reset, and account deletion. Check In with location. Gym ID and occupancy. Paid member and athlete registration. Test data for Objective 1 and the load test, with Pierre. Automatic accessibility checks and the flow tests. Project home page. Server requests to CTI.} \\
\addlinespace
\revised{Pierre G. Vales Torres} & \revised{Containers, database, Laravel API, and backend CI. Database tables and rules, and the fake test data. Staff and Administrator accounts: sign in, roles, invitations, password reset, and the audit log. Dashboard base: layout, menu, design tokens, and Spanish and English texts. Staff check-out and the automatic 2-hour close. Academic terms and the athlete roster, with the Word and PDF download. Staff panel: who is inside, assisted Check In, and visit corrections. Occupancy on the dashboard. Security updates and the review of merged pull requests.} \\
\addlinespace
\revised{Viviana V. Ramos Vivoni} & \revised{Wireframes in Figma for registration, Check In and Gym ID, sign in and account, and Staff screens; content, analytics, and administration screens are in progress. The open design decisions of the screen map. User app screens with Carlos: sign in, Gym ID, paid registration, emergency contact, and My visits. Manual accessibility check. Project website on GitHub Pages and the design document updates.} \\
\bottomrule
\end{tabularx}
\end{table}

\subsection{Testing Plan}

Each test belongs to one objective and to one or more requirement IDs (REQ-NF-44), and every case uses a controlled dataset. One seed command loads the same fake data into a fresh local database, and the automated tests and the live demo both run on it, so the demo shows the cases the tests prove.

\revised{\textbf{Unit and regression tests.} Unit tests check each part alone: Vitest for each screen, its states, and the shared code, and PHPUnit for each API route and database rule on PostgreSQL. Playwright browser tests [12] run each flow on the full stack. Continuous integration runs all of them on every pull request and every push to \texttt{main}, so each change reruns the tests of earlier work. The backend workflow also runs lint, migrations, seeding, a health check, and a check that the API answers 503 when the database is down.}

\textbf{Authorization (Phases 3 and 4):} a controlled dataset checks valid and expired paid memberships, valid and invalid athlete approvals, missing authorization, and unmatched users against expected approve and deny decisions. \textbf{Check In and sessions (Phase 3):} tests cover location inside and outside the boundary, denied permission, inadequate accuracy, successful, rejected, and duplicate Check Ins, Staff check-out, the automatic close after 2 hours, closures, and concurrent requests that must create at most one session. Staff-assisted Check In tests (Phase 5) verify the authorized user without a smartphone, photo and authorization checks, actor and timestamp audit, and no overlapping session.

\textbf{Occupancy (Phase 3):} expected active-session counts are compared after Check Ins, check-outs, automatic closes, Staff corrections, and concurrent events; User, Staff, and Administrator views must agree. \textbf{Roles (Phase 5):} the permission matrix confirms allowed actions succeed and forbidden actions fail. \textbf{Analytics (Phase 5):} each finalized metric is compared with a known result from controlled historical data. \textbf{Integration (Phase 6):} one integrated dataset runs from sign in to session creation to analytics, and all five areas are checked together. Engineering-constraint pass conditions and timing appear in Engineering Constraints and Verification.

\revised{\textbf{Concurrency (REQ-NF-42).} A backend test sends simultaneous Check Ins for one user, and repeated Authorization Number or roster claims, to the real database. It passes only with one session or one claim. The load tool cannot release all requests at the same instant, so this check is its own test.}

\revised{\textbf{Accessibility.} An axe-core scan [13] with the WCAG 2.0, 2.1, and 2.2 A and AA rules runs on each screen inside the browser tests and fails the build on a violation, and the linter's accessibility rules are on. A manual checklist covers what a scan cannot: keyboard only, visible focus never covered by the bottom bar, NVDA on Windows, VoiceOver on iPhone, 200\% zoom, a 320 pixel wide screen, and status messages read aloud. Each pull request that adds a screen includes the filled checklist. Scans find only part of the problems, so results claim no automatically detectable violations plus the manual checks, not full WCAG conformance.}

\revised{\textbf{Stress testing with k6 (REQ-NF-15, REQ-NF-42, REQ-NF-50).} k6 scripts [11] sign in the fake users before the timed part and then run three profiles. The gym gave no count for its busiest hour, so the team set the sizes from the estimate of about 150 visits a day: a peak of 50 Check Ins spread over 5 minutes (about a third of the daily visits in the busiest hour, compressed), a stress run of 100 over 5 minutes, and a burst of 50 members at once, each sending two Check Ins in parallel. In the peak and stress runs each member repeats the Check In 60 seconds later and must get the same session. A run passes with a Check In p95 under 3 seconds, errors under 1\%, and every check passing; after the run, the test counts sessions per user and reads the occupancy number. A preliminary run uses the local containers, where k6 and the app share one computer, so its numbers only compare runs with each other. The full run is on the CTI VM. Separately, the Check In button must be usable in under 3 seconds on mobile data on a mid-range phone (REQ-NF-16), which needs a real phone; the test phone and network are \pending{chosen by the team; all three members can test on iPhone and Android}.}

\revised{\textbf{Acceptance testing.} After the release is on the CTI VM, gym staff and real users, both paid members and athletes, run the main flows on their own phones and computers: registration, Check In, Gym ID, Staff-assisted Check In, check-out, and the store sync. Browsers: iPhone with Safari, Android with Chrome, and Windows with Chrome and Edge, with macOS and Firefox as a secondary check (REQ-NF-30, REQ-NF-43, REQ-NF-49). The gym records accepted or not accepted for each flow and the problems found, and each problem becomes an issue. Who takes part and when: \pending{participants and dates, waiting on the gym}.}

%=====================================================================
